% page 495 of the facsimile (image p-494 of the scan)
% block 175.3 x 259.3 mm ; left margin 26.7 mm
\pgc{5.080mm}{0.000mm}{
\l{}
\l{}
\l{}
\l{}
\l{\cel{54}487.}
\l{}
\l{\cel{3}\sou{regalio}. (eklezio katol., en Francia). Yuro quan havis la rejo}
\l{\cel{5}pri inkasar la revenui di la episkopeyi, di la abadeyi,}
\l{\cel{6}vakanta, e nominar pri la benefici qui dependis de li. -}
\l{\cel{6}DEFIRS.}
\l{}
\l{\cel{3}\sou{regardar}. (\sou{trans}.) Direktar sua okuli apertita (ad ulo od ulu).}
\l{\cel{6}EFI.}
\l{}
\l{\cel{3}\sou{regato}. Konkurso de bateli, qui avancas o per seglo, o per la}
\l{\cel{6}agado da remeri. - DEFIRS.}
\l{}
\l{\cel{3}\sou{regenerar}. (\sou{trans.)}\cel{2}I. Reproduktar (parto destruktita). - II.}
\l{\cel{6}Ri-novigar(+ri-nuvigar) etikale, morale. - DEFIRS.}
\l{}
\l{\cel{3}\sou{regento}. La persono qua guvernas dum la minoreso o dum la absent-}
\l{\cel{6}eso di la suvereno. - DEFIRS.}
\l{}
\l{\cel{3}\sou{regio}. I.Exploto, administro, di domeno ye la nomo di altru e ye}
\l{\cel{6}la profito da ica. - II. Imposto-inkaso agata nemediate da la}
\l{\cel{6}agenti di la stato, profitote da ica (kontraste kun la sistemo}
\l{\cel{6}qua konsistas ye luar la yuro inkasor la imposto-revenui, la}
\l{\cel{6}lokacinto pagante a la stato pekunio-quanto konvencionita, e}
\l{\cel{6}retenante por su la ecesajo. - III. La employataro, la kontori,}
\l{\cel{6}qui reletas ta percepto nemediata. - DEFIS.}
\l{}
\l{\cel{3}\sou{regimento}. (\sou{armeo)}. Trupo de plura batalioni od eskadroni, direkte}
\l{\cel{6}da kolonelo. - DEFIRS.}
\l{}
\l{\cel{3}\sou{regiono}. I. Parto plu o min granda di lando qua submisesas a kli-}
\l{\cel{6}mato komuna. - II. Parto determinita di korpo. (\sou{anke metaf.})}
\l{\cel{6}DEFIRS.}
\l{}
\l{\cel{3}\sou{registro}. Libro, kayero, sur la skriburo-virga pagini di qua onu}
\l{\cel{6}notas senfalie la fakti pri qui onu volas memorer. - DEFIRS.}
\l{}
\l{\cel{3}\sou{regnar}.\cel{2}(\sou{trans.}) Aplikar sua povo rejala, sua autoritato suve-}
\l{\cel{6}renala. - DEFIS.}
\l{}
\l{\cel{3}\sou{regolo}. (\sou{zool.)}\cel{2}Mikra pasero kantera, ek la familio \textquotedbl{}silvidei\textquotedbl{}.}
\l{\cel{6}-- L.\cel{1}\sou{troglodytes parvulus}.-}
\l{}
\l{\cel{3}\sou{regoliso.}\cel{2}(\sou{bot.)}\cel{2}Planto di qua la radiko uzesas kom dolcigivo}
\l{\cel{6}pri rumo, bronkito, e c. - L.\cel{1}\sou{glycyr-rhiza.}- IS.}
\l{}
\l{\cel{3}\sou{regresar.}\cel{1}-\cel{2}(\sou{netrans.}) Irar kontre progresado. - EFIRS.}
\l{}
\l{\cel{3}\sou{regretar}\cel{2}(\sou{trans.)}\cel{2}I. Chagrenar pro ne plus havar (ulu, ulo). -}
\l{\cel{6}II. Chagrenar pro agir (o pro ne agir) (ulo). - EF.}
\l{}
\l{\cel{3}regular. (\sou{trans.)}\cel{2}(\sou{tekn.)}\cel{2}Submisar ad ordino rigoroza (la fun}
\l{\cel{6}ciono di mekanismo, di mashino, e c.) - DEFIRS.}
\l{}
\l{\cel{4}regulatoro. (tekn.)\cel{2}Organo qua regulas la funciono di mekanismo,}
\l{\cel{6}di mashino. - DEFIS}
}
